\chapter{Se repérer et représenter l'espace}
\label{exercices:18}
\section*{Exercice 18.1 -- Coordonnées et perspective}
Dans un pavé repéré par $0\leq x\leq6$, $0\leq y\leq4$, $0\leq z\leq3$, situer $A(6;0;3)$ et $B(2;4;1)$. Les longueurs des arêtes fuyantes se mesurent-elles directement sur un dessin en perspective cavalière ?

\section*{Exercice 18.2 -- Section d'une boule}
Une boule de rayon $13\unite{cm}$ est coupée par un plan situé à $5\unite{cm}$ du centre. Calculer le rayon et l'aire de la section. Que se passe-t-il si la distance devient $13$ puis $14\unite{cm}$ ?

\section*{Exercice 18.3 -- Section et patron d'un cône}
Un cône de rayon $6\unite{cm}$ et de hauteur $8\unite{cm}$ est coupé à $2\unite{cm}$ du sommet par un plan parallèle à sa base. Calculer le rayon de la section et le rapport des volumes petit/grand. Calculer la génératrice et l'angle du secteur de son patron latéral complet.

\section*{Exercice 18.4 -- Sur la Terre sphérique}
Décrire la position d'un point de latitude $20\degre$ S et longitude $40\degre$ O. Avec une circonférence terrestre modélisée par $40000\unite{km}$, calculer la longueur d'un arc de $18\degre$ sur un grand cercle. La même différence de longitude donne-t-elle cette longueur sur tous les parallèles ?
